\chapter{Literature Review}

\section{Microwave Ablation in Tumor Treatment}

Microwave ablation employs electromagnetic waves (typically at 915\,MHz or 2.45\,GHz) to induce dielectric heating in biological tissue. Key treatment parameters include:

\begin{itemize}
    \item \textbf{Input Power} (W): Energy delivered to the tissue via the antenna.
    \item \textbf{Treatment Duration} (minutes): Time of energy application.
    \item \textbf{Antenna Design}: Determines the spatial distribution of the electromagnetic field.
    \item \textbf{Frequency}: Operating frequency of the microwave system.
\end{itemize}

The resulting ablation zone is characterized by its \textbf{diameter} (or width), \textbf{length} (longitudinal extent), and overall \textbf{shape} (spherical vs.\ ellipsoidal), which is quantified by the \textbf{sphericity index} (diameter/length ratio).

Several antenna designs have been developed for MWA, including monopole, dipole, dual-slot, triaxial, floating-sleeve, and helical dipole antennas. Each design produces distinct ablation zone geometries based on the specific absorption rate (SAR) pattern it creates.


\section{Machine Learning in Medical Treatment Planning}

Machine learning has been increasingly applied to medical treatment planning, including:

\begin{itemize}
    \item Radiation therapy dose prediction using deep learning.
    \item Surgical outcome prediction using ensemble methods.
    \item Drug response prediction using random forests and neural networks.
    \item Temperature prediction during thermal ablation using deep learning (e.g., the Triple Coaxial-Half-Slot antenna study, 2021) \cite{triple_coaxial_2021}.
\end{itemize}

For structured tabular data with limited sample sizes, \textbf{tree-based ensemble methods} (Random Forest, Gradient Boosting) consistently outperform deep learning approaches. This is supported by extensive benchmarking studies and is particularly relevant to our dataset.


\section{Existing Approaches and Research Gap}

Most existing ablation zone prediction methods rely on:

\begin{enumerate}
    \item \textbf{Physics-based simulations} (FEM/FDTD): Accurate but computationally expensive (hours per case).
    \item \textbf{Manufacturer lookup tables}: Quick but limited to specific antenna--power--time combinations.
    \item \textbf{Empirical formulas}: Simple but poor generalization across antenna designs.
\end{enumerate}

\textbf{Research Gap:} There is no comprehensive ML-based predictive model that generalizes across multiple antenna types, power levels, and treatment durations using aggregated data from multiple published studies. This project addresses this gap.
